\section{上下文与成本（RQ1）}
\label{sec:context}

让每个智能体的提示词只追加，使每次运行的成本降低了三分之一到一半，得分不变。之后测量的上下文手段让成本变化了
7\% 到 20\%，没有一个改变准确率（表~\ref{tab:levers}）。

\para{只追加的提示词}旧编译器每一轮都重新渲染指令，因此 30\% 到 40\% 的请求在系统部分就断开了缓存。在 9 月 30 日的
一次配对运行中（同一次重启、同样负载，20 次运行，全部解出），纪元方式有 89\% 到 94\% 的提示词 token 读自缓存，旧方式
为 76\% 到 81\%。旧方式每次运行的成本高出 50\% 到 108\%，例如一道单人 SQL 题是 0.81 美元对 0.44 美元。纪元边界之外的
前缀改写只占 9,897 次请求的 0.25\%；空闲长达一小时之后，缓存命中率仍有 98\% 到 99\%。

\para{服务商的缓存需要正确的请求头}在 ChatGPT 后端上，只带缓存键什么也缓存不了。加上会话和线程请求头之后，命中率
升到 97.8\% 到 99.6\%；重放轮次状态又让一轮中途的请求从 0 升到 98.5\%。在这些修复之前，一个长期运行的团队一周花掉了
5.55 亿个提示词 token：它在 872k token 窗口的 85\% 处才压缩，而后端并不支持这么大的窗口；它的提示词涨到 744k token，
而且完全没有命中缓存。现在压缩在 231k token 时触发。

\para{委派的代价在冷启动}在第 2 批（9 月 29 日）中，委派的运行快了 25\%，但发送的提示词 token 是原来的 2.46 倍，
未命中缓存的 token 是 7.57 倍。52 个被派生的智能体，每个的第一次请求都约有 42k token 未命中缓存，占这些运行未命中
token 的 57\%。随后做了两项修复：智能体自己的笔记改为放在共享指令之后、作为用户消息发送（一个审查者第一次请求中
命中缓存的 token 从 42k 里的 18k 升到 38k）；克隆体继承源智能体的工具和纪元（审查者第一次请求的缓存命中率从 13\% 升到
96\%）。再加上团队共用的缓存键、静态的交接工具和一个预热关口，每次委派运行的未命中 token 从 643k 降到 390k。

\para{延迟跟着 token 走，不论是否命中缓存}在同一次运行、同一时刻发出的 2,354 对匹配请求上，首个事件的到达时间每多
10k 个命中缓存的 token 增长到 1.09 倍，每多 10k 个未命中的 token 增长到 1.11 倍；消息条数没有影响。领队是关键路径：
一次运行约 90\% 的活跃时间花在它自己的模型调用上，它 86\% 的轮次只调用一个工具。减少领队的轮次能省时间；缩小
提示词主要省钱。

\begin{table}[t]
\centering
\footnotesize
\caption{V1 任务集上的上下文手段。比值为实验组除以同批基线；“完全解出”指拿到满分的运行。}
\label{tab:levers}
\begin{tabularx}{\textwidth}{@{}>{\raggedright\arraybackslash}p{36mm} >{\raggedright\arraybackslash}p{18mm} L L l@{}}
\toprule
手段 & 批次与规模 & 成本或 token（对基线） & 得分与时间 & 结论 \\
\midrule
只追加的纪元（对旧编译器） & 9/30，20 次运行 & 成本低 33\% 到 52\% & 两边全部解出 & 默认开启 \\
提前写项目笔记（播种运行） & b101，5 道题 & token 0.84 & 时间 0.83；完全解出 5 对 4 & 采用 \\
读取指引（不整文件加载） & b96，24 对 & 未命中 0.80（$p = 0.001$），成本 0.86 & 时间 0.95；完全解出 18 对 16 &
  采用 \\
从未调用的工具只给简短 schema & b97，24 对 & 领队首次请求 28.9k 降到 21.9k token；token 0.89，成本 0.93 & 时间 0.97；
  完全解出 18 对 15 & 采用 \\
精简工具结果 & b100，24 对 & token 0.98，成本 0.98 & 时间 1.06 & 放弃 \\
多个任务共用一个会话 & s1，5 道题 & 成本 1.06 & 得分相同 & 放弃 \\
60k token 时移出旧工具结果 & b95，24 对 & 未命中 1.29，成本 1.14 & 时间 1.02 & 放弃 \\
\bottomrule
\end{tabularx}
\end{table}

\para{摘要会丢规则}对 861 份已提交摘要的审计发现，摘要长度中位数只有原文的 1.1\%，规则句的 token 只有 36.9\% 被原样保留；
每一份后续摘要大约只保住前一份所保留内容的一半。现在摘要器会原样保留当前任务、对任务的修改和最新的汇报。对话回忆
测试显示不出这项改动的收益：每个实验组在三次摘要之后都记住了全部八条植入的事实，所以这项改动的依据只有那次审计。

\para{RQ1 的回答}上下文管理是一个成本手段。让提示词前缀稳定，比之后所有削减 token 的做法加起来都值；而一个打断
缓存的削减，花掉的比省下的多（第~\ref{sec:discussion}~节）。至于时间，要紧的是领队的轮次，而不是提示词的大小。
